% 0_abstract.tex
\chapter*{Abstract}

Procurement runs on documents, and every defective one that slips through is far more expensive to fix downstream than to catch on receipt. This dissertation takes on one such gate at Hilti AG: validating each Advanced Shipping Notice (ASN) against its originating Purchase Order (PO). Today that check is a binary accept or reject filter on static fields; it performs no semantic and no cross-document validation, and when it fails it draws a buyer into a manual reconciliation loop of 30 to 90 minutes per defect.

Large language models are tempting here, since they read supplier specific fields and judge contract terms with ease, but they are unreliable on the arithmetic the task also demands. Recent work couples them with deterministic verification, yet no study has compared the major hybrid designs head to head, leaving practitioners without the accuracy, cost and latency trade-offs that decide which to build. This dissertation closes that gap. It formalises the \emph{enforcer pattern}, a deterministic layer that holds authority over the verdict of every rule a computation can settle, and delivers the first reported benchmark of the major hybrid validation designs on a single 22 rule task.

Framed as design science research, the artefact is a 7 stage pipeline that replaces the binary gate with a tiered PASS, REVIEW or REJECT decision. It is evaluated under 5 designs that differ in whether the enforcer is present and where it acts, on 67 deterministically generated synthetic cases and 32 real cases whose ground truth is the consensus of 3 procurement analysts. A single language model is held fixed across all designs, and a second model from a different family is used to test robustness, with the designs compared on disposition agreement, per rule accuracy, cost and latency under bootstrap confidence intervals.

The enforcer separates the designs. On the 32 real cases the 3 designs that carry it reach 0.88 disposition accuracy against 0.81 without it, a directional gap on this small sample; on the 67 synthetic cases, where the contrast is statistically reliable, that gap widens to 20 points. Where the enforcer sits is an operational choice rather than an accuracy one: settling the deterministic rules before the model call gives the single optimal design, matching the best accuracy and recovering the most rules while leading on cost, at EUR 0.073 per ASN, and latency, at 7.89 seconds median. A supplier history prior adds nothing once the enforcer is present, and under the second model the enforcer's advantage does not merely survive but widens, the clearest indication that the gain comes from the architecture and not the model. The work contributes a reusable architectural pattern, the first reported benchmark of the major hybrid validation designs on a single 22 rule task, and an evidence based deployment recommendation for procurement document validation.

\vspace{2em}
\noindent\textbf{Keywords:} large language models; hybrid LLM and deterministic pipeline; enforcer pattern; document validation; procurement; design science research.
